\chapter{Proposed Research}
\label{proposedRes}

To date, we have developed two formulations for the LRS problem (the first formulations
as far as we know), identified valid inequalities that strengthen the formulations, and
developed a branch-and-price algorithm to solve instances. We have demonstrated that
the algorithm can solve to optimality a variety of instances up to 30 customers 
and 5 facilities. In this section, we describe our proposed research. 

In section \ref{algo-dev}, we describe our plans for improving the implementation
of the branch-and-price algorithm to solve larger instances. In section 
\ref{sepecial-cases}, we present special cases of the LRS problem and explain how
we can modify our current work to solve these special cases. In section 
\ref{extentions}, we  present our ideas for the extensions of the LRS problem.   

\section {Algorithmic Improvements}
\label{algo-dev}
We plan to further develop our branch-and-price algorithm with the goal of speeding 
up our algorithm and solving larger instances. To do so, we will  
improve our column generation algorithm, the methods for obtaining upper bounds 
at the nodes of the tree and choosing a node to process in the branch-and-price tree.

\begin{enumerate}
\item {\bf Column generation} The branch-and-price algorithm spends the largest amount
  of time on the column generation algorithm at each node of the tree. \citet{Vanderbeck2}
  lists several reasons why this may happen.
  One reason is that, as the solution gets closer to 
  optimality, it becomes more difficult to find columns with negative reduced costs; 
  this is known in the literature as the tailing-off effect. Another reason is the 
  degeneracy in the RMP solution. The pairing structure causes degenerate solutions.
  There may be two sets of pairings visiting the same customers with the same total costs  
  but including different pairings. Degeneracy in the primal solution results in 
  multiple optimal dual solutions. Therefore, for several iterations,  new 
  columns are generated using different dual solutions, but the optimal cost of the 
  RMP may not improve because the optimal primal solution is moving from one degenerate 
  solution to another. To decrease the time spent on the column generation algorithm, 
  we will make the following improvements to our current algorithm.

  {\bf Improving the Initial Set of Columns:}

  We will try to improve the initial set of columns that are provided to 
  initialize the RMP at the root node. A better set of initial columns may result 
  in better dual variables that can create columns that can be a part of the optimal
  RMP solution. Therefore, a better initial set may reduce the number of iterations 
  of the column generation algorithm. Currently, we are 
  generating all feasible columns that include one customer, two customers and 
  three customers. Each column represents a path in which the vehicle does not 
  return to the depot before serving all customers in the path. However, since 
  in our LRS problem, a column (pairing) may represent a set of routes, we may 
  be able to improve the initial solution by generating columns including more 
  than one route. This may result in needing fewer vehicles to serve all of the 
  customers. Our idea is to use a heuristic designed for bin packing problems 
  to combine the initial routes with one, two and three customers to create 
  columns including more than one route.  
    %The column generation algorithm starts with some initial set of columns and 
    %creates columns based on the dual variable information. In the 
    %first iterations, the quality of the created columns depends on the quality of 
    %the initial columns. Therefore, if we generate better initial columns then
    %the algorithm will skip some of the   
  

  {\bf Improving the Management of Columns:}
  \begin{itemize}
  \item Currently, we are solving the pricing problem to optimality during each 
    iteration of the column generation algorithm. That means that we are finding 
    the column with the most negative reduced cost as well as all possible columns 
    with negative reduced costs. Although the column with the most negative reduced 
    cost potentially reduces the objective function value by the largest value,  
    solving the pricing problem to optimality gets more difficult as the number of 
    customers in the problem increases. 
    %We tried terminating the pricing problem after some number of columns generated. This 
    %premature exit of the pricing problem did not give better solution times even 
    %we saved time at each pricing problem. 
    Our idea is to implement an approach similar to \citet{Larsen}. He creates a policy 
    in which he sets two parameters, $m$ and $n$. He solves the pricing  problem to 
    optimality until $m$ columns in total are added to the RMP, then he exits the pricing 
    problem when optimality is reached or at least $n$ columns are generated. Therefore, 
    he takes advantage of the best columns as well as the advantage of early termination. 
    The parameter values need to be determined empirically. 
    %The maximum number of 
    %columns to be created before allowing premature exit of pricing problem 
    %should be determined empirically.   

  \item At each step of the column generation algorithm, typically more than one column is
    generated, so more than one column typically is added to the RMP at each step.
    The hope is that adding multiple columns will improve the solution 
    of the RMP and reduce the number of iterations of the column generation algorithm.
    However, since some of the columns might be similar to each other,  
    adding all of them might be redundant. Furthermore, since solving the RMP becomes
    more difficult as the number of columns increases, it is important
    to add the right number of columns at each step of the column generation
    algorithm. We will experiment with different policies for adding columns, 
    such as adding the best column, adding the last $n$ columns in the list and 
    adding $n$ random columns with negative reduced cost, and we will
    compare the performance of the algorithm with different values of $n$.  

    %Solving the pricing problem to optimality finds all possible columns 
    %with negative reduced cost and therefore finds the column with the most 
    %negative reduced cost that reduces the objective function most. However, 
    %as the number of customers in the problem increases, solving the pricing 
    %problem to optimality gets more difficult and therefore it takes more time
    %to get the best column. Hoping to improve the solution of the RMP and to
    %skip some iterations of the column generation algorithm, adding more than 
    %one column at each step of the column generation algorithm is suggested in
    %the literature. However, some of the columns might be similar to each other, 
    %therefore adding all of them  might be redundant.   
  \item In many column generation implementations, researchers report that 
    creating a column pool decreases computation time. A column pool stores
    columns that are generated but not added to the model. The idea is that 
    instead of solving the pricing problem at each iteration, we first search 
    the column pool for the columns with negative reduced costs calculated using 
    the current dual variables. If there are columns, then we will select some 
    of those to add to the model. Otherwise, we will solve the pricing problem
    to get columns with negative reduced cost. 
   
   
   %\item Creating column pool is one of the most common implementation improvements in branch-and-price algorithms.
   % Currently we have a column pool, but this pool only includes the columns that are added to the formulation. 
   % We are solving our pricing problem that is an ESPPRC to optimality and get the column that have most negative
   % reduced cost and discard other feasible columns. We will update our column pool and store some of the columns 
   % that are created in pricing problem and have negative reduced cost to use in next column generation steps. Also,
   % we can solve the pricing problem for one facility and can use the column pool for the other facilities. 
   % Instead of solving the ESPPRC for other facilities, we can simply check the reduced costs of columns and add 
   % these columns to the formulation. 
%**********************************888
  %\item For large instances, we will not always solve the exact ESPPRC. We may terminate the algorithm as soon 
   % as we find a column with negative cost. Also, if the pricing problem takes too much time for large instances 
   % we will develop a heuristic to find feasile columns. In this heuristic, we might relax one of the resource 
    %constraints and solve the ESPPRC. 
 %****************************************8   
   % We have a path pool in which we store the paths that are included in at least one 
   % pairing. We can constuct proper sets of paths in which the paths can be combined to obtain pairings. Let 
   % P1:(source-1-2-3-sink) and P2:(source-3-4-sink) be two paths. Since node 3 is in both paths, they cannot
   % be combined to get a feasible pairing. We can solve knapsack problems over these proper sets to get a 
   % pairing with negative reduced cost.   
 %***********************************************   
  \end{itemize}

\item {\bf Upper bound (incumbent solution)} In branch-and-price algorithms,
  having a good upper bound (incumbent solution) for the problem is very 
  important since it can help to decrease the size of the tree. We are currently 
  solving the IP model at the root node to optimality to find an upper bound, and
  using the default primal heuristic in MINTO to find an incumbent solution at
  the other nodes of the tree. One idea is to solve the IP models at the nodes 
  to optimality from time to time to improve the upper bound. Another idea is to 
  develop an application specific heuristic to find an upper bound. We can do 
  this using VRP solvers to create routes in combination with bin packing 
  heuristics to combine the routes.
  %Our idea is to develop an application specific heuristic to 
  %find an upper bound. We can do this using VRP heuristics to create routes
  %in combination with bin packing heuristics to combine the routes. 
\item {\bf Node selection policy} The order in which the nodes of the 
  branch-and-price tree are processed is one factor that may impact the 
  efficiency of the algorithm. Currently, we are using a best-bound node selection
  policy, which seeks to improve the lower bound of the problem. Another common
  selection policy is depth-first, which seeks to find an upper bound 
  (incumbent integer solution) as quickly as possible in hopes of eliminating 
  some of the branches of the tree. In some limited tests of using the depth-first 
  search policy, we obtained worse computation times. However, it may be 
  beneficial to use a combination of these two policies as in \citet{Vanderbeck}.
  He first applies the depth-first search to get an incumbent solution. He uses 
  the depth-first policy until the gap between the incumbent solution and the best 
  bound in the tree is less than some prespecified value. After, he switches to 
  the best-bound policy. 
\end{enumerate}
By implementing these improvements, we hope to be able to solve instances
with at least 50 customers and 10 facilities in a reasonable time.   

\section{Special Cases of the LRS Problem}
\label{sepecial-cases}
The LRS problem generalizes a number of problems that appear in the literature. 
By relaxing some of the constraints and fixing some of the parameters of the 
LRS problem, we can obtain some special cases, such as location-routing problems 
(LRP), multi-depot vehicle routing problems (MDVRP) and vehicle routing problems with 
multiple trips (VRPMT). For some of these problems, our algorithm can provide
upper bounds, lower bounds or exact solutions, which would improve the results
reported in the literature and contribute to the performance evaluations. In the 
following sections, we explain the necessary modifications  to obtain those particular 
problems and describe our potential contributions.

\subsection{LRP}

There are few exact algorithms for LRPs, especially for those with 
capacitated facilities, capacitated vehicles and symmetric distance matrices, 
and the literature only reports exact solutions for small instances. Therefore, 
an algorithm that solves larger instances of capacitated LRPs to optimality 
will add great value to the literature. \citet{Laporte86} report solving 
to optimality instances of symmetric LRPs (with capacitated facilities
and capacitated vehicles) with at most 20 customers and 8 depots. \citet{Hansen} 
give optimal solutions for instances with 15 customers and 3 depots. 
Clearly, there is a need for algorithms that can solve LRPs of practical 
size, which is why many papers develop heuristic algorithms. Because exact algorithms
are not available for large problems, however, the results of heuristics 
often are evaluated only by comparing them with the results of other heuristic algorithms. 
%because LBs are not availale. 
\citet{Perl} present lower bounds for their instances (up to 55 customers 
and 15 depots) and report average 5.2\% gap between their heuristic results 
and lower bounds. However, \citet{Srivastava}, \citet{Srivastava2} and 
\citet{Tuzun} compare their results with other heuristics and do not report 
lower bounds. Therefore, being able to obtain lower bounds and 
upper bounds (if  not optimal solutions) from our algorithm would allow researchers to 
better evaluate 
the performance of their heuristics. 

To construct an LRP from the LRS problem, we drop the time constraints for 
the vehicles and assume a one-to-one relationship between vehicles and routes 
in the relaxed model. For an {\em LRP with capacitated facilities and capacitated 
vehicles}, we can use the set partitioning-based formulation AUG-SPP-LRS 
but with a different interpretation of a column. A column will represent a
feasible vehicle route instead of a pairing (a set of feasible vehicle routes). 
A vehicle route is feasible if the vehicle starts and ends at the same depot, 
the sum of the demands of the customers on the route is at most the vehicle 
capacity and each customer is visited at most once. Changing the column
definition changes the pricing problem, so the network defined in section 
\ref{form-pricing-prob} must be modified to create columns representing 
the vehicle routes. The modified pricing problem will still be an elementary 
shortest path problem with resource constraints,  but the vehicle capacity 
will be the only resource. In the network, all arcs leaving the sink node 
are deleted, which means we will not allow the sink node to be visited more 
than once. In  this case, a path originating at the source must end at the 
sink but cannot be extended from the sink node. The modified pricing problem 
will be easier to solve than the current pricing problem, since the number of 
resources is smaller and the number of possible paths that can be created 
is smaller. 

Furthermore, by removing constraints (\ref{spp_con2}) and (\ref{lb-open-fac}) 
from the AUG-SPP-LRS formulation, we can get an 
{\em LRP problem with uncapacitated facilities}. Dropping the constraints 
requires that we update the definition of the reduced cost of a column. 
The reduced cost will be similar to the expression (\ref{red-cost}), except
that it will not include the dual variables associated with the dropped 
constraints: 
$\hat{C}_{jp}=C_{jp}-\sum_{i\in N}a_{ipj}.\pi_{i}+\sum_{i \in N}a_{ipj}.\sigma_{ji}$.

We will modify our current model and solution algorithm to solve capacitated and
uncapacitated LRPs. We expect to be able to solve or at least provide a feasible
solution  for instances
as large as 100 customers. 
%We will provide lower bounds for the location
%routing benchmark problems and compare our optimal or best integer solutions 
%with the solutions provided so far. 

\subsection{MDVRP}

As noted in the literature survey section for the VRP (section \ref{lit-vrp}), 
there are only two papers that develop exact algorithms for the MDVRP. 
\citet{Laporte6} solve symmetric MDVRPs with instances as large as 
 50 customers and 8 depots using branch-and-bound algorithm. \citet{Laporte7} 
solve asymmetric MDVRPs with instances up to 80 customers and 3 depots. 
As with the LRP, heuristics are used to solve large instances (up to 360 customers 
and 9 depots) and they are evaluated by comparing their results with the best 
reported results. Because lower bounds for the instances are not known, the 
integrality gaps are not known. Using our formulation and solution algorithm, 
we hope to be able to generate at least good lower bounds if not optimal
solutions for MDVRPs. Because set partitioning-based models tend to generate 
tight lower bounds, even generating lower bounds would provide another measure 
for the quality of the existing heuristic solutions. 
%answer the question of how good a heuristic solution is compared to the optimal solution.        

To define an MDVRP from the LRS problem, we drop the time constraints for the 
vehicles. In addition, we set the fixed cost of the facilities to zero.  Moreover, 
in an MDVRP, it is assumed that the facilities are already constructed. Therefore, 
we no longer need the location variable $T_j$ that shows whether the candidate 
location $j$ is selected as a facility or not. We can simplify our set 
partitioning-based formulation (AUG-SPP-LRS) to construct a formulation for the MDVRP
by removing the $T_j$ variables and the constraints (\ref{lb-open-fac}) and (\ref{relation}). 
Similar to the LRP, a column in the MDVRP formulation represents a feasible vehicle 
route. Therefore, we will make the same modifications to the pricing problem to 
create feasible vehicle routes instead of pairings. Here, however, the reduced cost 
definition of a column will change to 
$\hat{C}_{jp}=C_{jp}-\sum_{i\in N}a_{ipj}.\pi_{i}+\sum_{i \in N}a_{ipj}.d_{i}.\mu_{j}-v_j$
in which we do not have the dual variable $\sigma_{ji}$ associated with the linking 
constraint (\ref{relation}). 
%Since in an MDVRP, it is assumed that the facilities are already constructed, we will no 
%longer need the location variable, $T_j$ that shows whether the candidate location $j$ is 
%selected as a facility or not. The set partitioning formulation that we will use to solve 
%an MDVRP is:
\begin{optprog}
{\bf(SPP-MDVRP)} Minimize & \objective{\sum_{j \in M} \sum_{p\in P_{j}} C_{jp} Z_{jp}} \label{mdvrp-obj}\\
subject to & \sum_{j \in M} \sum_{p \in P_{j}} a_{ipj} Z_{jp} & = & 1 & \forall i \in N \label{spp_mdvrp} \\
           & \sum_{p \in P_j} \sum_{i \in N} a_{ipj} d_i Z_{jp} & \leq & CapF_{j} & \forall j \in M \label{spp_mdvrp2} \\
           & \sum_{p \in P_j} Z_{jp} & = & nV_j & \forall j \in M \label{tot-veh-mdvrp} \\
           & L_j \leq nV_j & \leq &  U_j & \forall j \in M \label{veh-bounds-mdvrp} \\
           & Z_{jp} &  \in & \{0,1\} & \forall j \in M, \forall p \in P_{j} \label{bin-mdvrp} \\
           & nV_j & \in & \mathbb{Z}^{+} & \forall j \in M \label{int-mdvrp}
\end{optprog} 
In constraints (\ref {veh-bounds-mdvrp}), $L_j$ is the minimum number of vehicles and 
$U_j$ is the maximum number of vehicles 
at facility $j$. Recall that we introduce the variable $nV_j$ to branch on the number 
of vehicles in section \ref{branchingRule} (under branching rule 2). The constraints 
(\ref{veh-bounds-mdvrp}) force lower and upper bound for the number of vehicles at 
each facility. The definitions of the other constraints are the same as in the SPP-LRS. 
 
We will use the modified pricing problem to solve the model SPP-MDVRP. We expect to 
be able to solve instances up to 100 customers. For larger problems, we may provide 
lower bounds and/or best integer solutions.  

\subsection{VRPMT}

The vehicle routing problem with multiple trips is another special case of the LRS
problem. Compared to the VRP, the VRPMT has received little attention in the 
literature. There are only a few recent papers and they develop heuristic algorithms 
for the VRPMT (see section \ref{lit-vrps}). All of the papers use the benchmark problems 
created by \citet{Taillard} that include between 50 to 199 customers. These problems 
are solved with a fixed number of vehicles and a predefined travel time limit for the 
vehicles. For feasible problems, researchers compare their heuristic solutions with 
the results of the existing heuristics based on a measure calculated using the 
ratio of the total travel length to the best VRP solution. Some instances are 
infeasible for a given number of vehicles and time limit. For those, they 
allow overtime with an associated penalty to find vehicle trips and they compare the 
length of the longest trip. 
%With our formulation and solution algorithm, we may be
%able to calculate the optimal number of vehicles and the optimal total travel time 
%for some of these benchmark problems.   
 
The LRS problem with one facility can be transformed into a VRPMT by fixing the 
location of the facility and setting the fixed cost of the facility to zero. 
Since there is only one facility and no fixed cost, we remove the location variables 
and all occurrences of the index $j$ (that corresponds to a facility) from the 
SPP-LRS formulation. Then, using our original pricing problem that creates feasible 
pairings, we can use the following set partitioning formulation as the master problem.
\begin{optprog}
{\bf(SPP-VRPMT)} Minimize & \objective{\sum_{p\in P} C_{p} Z_{p}} \label{vrpmt-obj}\\
subject to & \sum_{p \in P} a_{ip} Z_{p} & = & 1 & \forall i \in N \label{spp_vrpmt} \\
           & \sum_{p \in P} Z_{p} & \leq & U \label{bound-veh-vrpmt} \\
           & Z_{p} &  \in & \{0,1\} & \forall p \in P \label{bin-vrpmt} 
\end{optprog} 
where $U$ is the number of vehicles available at the depot. In the SPP-VRPMT, there 
is no capacity constraint for the facility since we are assuming that the capacity 
of the facility is enough to serve all customers. 

\citet{Olivera} provide a different set partitioning-based formulation of the VRPMT, 
in which the decision is how to group routes into pairings. In their formulation,
the columns represent routes and they include a constraint to group these routes 
in order to generate a feasible task to assign one vehicle. In our pricing problem, 
we will create pairings, each of which represents a feasible assignment for one vehicle.

We will use the model SPP-VRPMT and our current column generation algorithm 
to solve the benchmark instances of VRPMT. For instances up to 100 customers, we 
will try to calculate the optimal number 
of vehicles and the optimal total travel time and evaluate the performance
of our algorithm compared to others in the VRPMT literature.
% We will try to provide optimal
%solutions or best integer solutions and lower bounds and evaluate the performance
%of our algorithm compared to others in the VRPMT literature.

\section{Extensions to the LRS problem}
\label{extentions}

The LRS problem can be extended to other problems by changing some of the 
assumptions and specifications of the problem. 
First, we will consider a version of LRS with stochastic demands. 
We will investigate how to integrate the stochastic demands to our set partitioning 
model and the methods to solve the stochastic LRS problem. 
Second, we will relax the 
assumption of identical vehicles. We will solve the LRS problem with 
a heterogeneous fleet. By differentiating the pairings based on the vehicle
type, we can use our current formulation. We will solve different pricing problems
for each vehicle type. 
Third, we will investigate the VRP with full truckloads. Modifying our
pricing problem and set partitioning model, we will provide an algorithm to solve 
the VRP with full truckloads. 

% extension that we
%want to work on is the LRS problem in which we do not assume ideantical vehicles. 
%Based on the related literature, the most common extensions for the 
%related problems are obtained by adding stochasticity to the model parameters. We
%will model a stochastic LRS problem as our second extension. Finally, we will investigate 
%the extension of our problem to the vehicle routing problem with full truckloads. 
%The following sections briefly describe the necessary changes and updates in our 
%current model and algorithm to obtain the extended problems.  

\subsection{Stochastic LRS Problem}

The current LRS problem assumes that all of the parameters are deterministic. 
We will investigate extending our current model to include stochastic parameters,
in particular, stochastic customer demands. Models with stochastic parameters are 
more difficult to solve than the corresponding models with deterministic 
parameters. For example, there are fewer papers studying stochastic versions
of the LRP than papers studying deterministic versions of the LRP (see \citet{Min}). 
Most of the stochastic papers consider just a single uncapacitated 
facility because of the difficulty of the problem. We may also consider relaxing 
some of the constraints in our stochastic model. 
  
To incorporate stochastic demands, our idea is to extend our set partitioning 
formulation to a stochastic programming model. \citet{Novoa} uses a 
set partitioning model to provide heuristic solutions for the VRP with
stochastic demands. We will investigate how to use column generation algorithm and
whether we can use our current pricing problem or not.  
We will develop an appropriate recourse policy to handle failures that occur in the constructed routes due to 
insufficient vehicle capacity.  
%\citet{Novak} use a set partitioning 
%model to solve stochastic VRP heuristically.

% We will determine whether we can use our current
%pricing model or how we can incorporate stochastic demands to our pricing problem. 
%We will investigate the ways of applying column generation for a stochastic master
%problem. 

  

\subsection{LRS Problem with Heterogeneous Vehicles}
The current LRS problem assumes that the vehicles are identical in capacity, 
time limit and fixed and variable costs.  However, the set partitioning-based 
formulation and the structure of the pricing problem can be extended to model 
non-identical vehicles, in particular, a heterogeneous set $S={1, ..., K}$ of 
vehicles, each with different properties. A vehicle of type $k$ has a capacity 
value $CapV_k$, a fixed cost $VF_k$ and an operating cost $OC_k$. In addition, 
the different types of vehicles may have different travel times on an arc as 
well as different working hour limits. Therefore, the travel time of a vehicle
of type $k$ from location $i$ to location $j$ is $t_{ijk}$ and the time limit 
for a vehicle of type $k$ is $TL_k$. 
 
We redefine our decision variables and parameters to incorporate the 
non-identical vehicle types. 
\begin{eqnarray*}
Z^k_{jp} & = & \left \{ \begin{array}{ll}
                 1  & \mbox{if pairing p of vehicle type $k$ is selected for facility j, $\forall p \in P^k_{j}$, $k \in S$ and  $j \in M$ }\\
                 0  & \mbox{otherwise}                                                  
                \end{array}
       \right. \\
a_{ipj}^k & = & \left \{  \begin{array}{ll}
                 1 & \mbox{if demand node $i$ is in pairing $p$ of vehicle type $k$ located at facility $j$, $\forall i \in N,$} \\
		 \mbox{}  & \mbox{$ j \in M, p \in P^k_{j}$ and $k \in S$ } \\
                 0 & \mbox{otherwise}
                         \end{array}
              \right. \\
C^k_{jp} & = & \mbox{cost of pairing $p$ of vehicle type $k$ associated with facility $j$, $\forall j \in M, p \in P^k_{j}$ and  $k \in S$ }
\end{eqnarray*}
Then, the set partitioning model for the LRS problem with heterogeneous vehicles will be:
\begin{optprog}
{\bf(SPP-LRS-HE)} Minimize & \objective{\sum_{j \in M} FC_j T_{j} + \sum_{j \in M} \sum_{k \in S}\sum_{p\in P^k_{j}} C^k_{jp} Z^k_{jp}} \label{he-obj}\\
subject to & \sum_{j \in M} \sum_{k \in S}\sum_{p \in P^k_{j}} a_{ipj}^k Z^k_{jp} & = & 1 & \forall i \in N \label{he_con1} \\
           &  \sum_{k \in S}\sum_{p \in P^k_j}\sum_{i \in N} a_{ipj}^k d_i Z^k_{jp} & \leq & CapF_{j} T_{j} & \forall j \in M \label{he_con2} \\
           & \sum_{j \in M} T_{j} & \geq & nREQ \\
           & \sum_{k \in S}\sum_{p \in P^k_{j}}a_{ipj}^k Z^k_{jp} &\leq & T_{j} &\forall j \in M, \forall i \in N \\ 
           & Z^k_{jp} &  \in & \{0,1\} & \forall j \in M, \forall p \in P_{j}, \forall k \in S \label{he-binary1} \\
           & T_{j}  & \in & \{0,1\} & \forall  j \in M \label{he-binary2}
\end{optprog}   

The interpretations of the constraints in SPP-LRS-HE are same as the 
corresponding constraints in the AUG-SPP-LRS, except that the type of the vehicle
is differentiated by an index $k$. The pricing problem to create columns for the 
SPP-LRS-HE will be the same as our current pricing problem except that we will solve 
one pricing problem for each type of vehicle as well as for each facility. 

If we assume that the travel times and the operating costs are the same for each 
type of vehicle, e.g., $t_{ijk}=t_{ij}$ and $OC_k=OC$, and there is no upper 
bound on the available number of vehicles of each type at each depot, then 
we can use our original set partitioning model AUG-SPP-LRS with the decision 
variable $Z_{jp}$ instead of using a different variable for each type of vehicle 
($Z^k_{jp}$). For a generated pairing, we will choose the cheapest feasible vehicle 
type across all vehicle types. Therefore, the type of the vehicle and the routes 
in the pairing will be determined in the pricing problem. The master problem will 
only be aware of the cost of the pairing but not be able to differentiate the 
type of the vehicle used for the pairing.  



\subsection{Vehicle Routing Problem with Full Truckloads}

The VRP with full truckloads (VRPFT) is a variation of the VRP that is studied 
in a few papers, e.g., \citet{Desrosiers2}, \citet{Arunapuram} and \citet{Gronalt}. 
The VRPFT finds the minimum cost set of vehicle routes, each of which starts and ends at 
the same depot and carries a given number of full truckloads between specified origin 
and destination locations. Since each truckload between each origin and destination 
locations requires one vehicle, the problem actually minimizes the 
total travel cost of empty vehicles. The VRPFT is similar to the LRS problem with
uncapacitated vehicles and without the restriction that each customer must be 
visited once. A customer must be visited in the number of truckloads that must picked up
from the customer. A network in which each arc between an origin and a destination   
location is represented with one node can be constructed. In this network, the nodes 
can be visited more than once to be able to satisfy the given number of truckloads. 
We will investigate the differences between the solution methods in the literature 
and a solution algorithm that we will obtain by modifying the current set 
partitioning model and the pricing problem.       



%\begin {itemize}
%\item We will construct a problem in which a path can include at most two customers. With this
%restriction, ESPPRC will be easier since possible number of pairings will decrease. This modification
%may help solving large instances.
%\item We will search the solution methods for LRS problem if we have stochastic parameters such as 
%stochastic demand or stochastic travel times. 
%\item We will solve the problem for two different type of products, let prodA and prodB be the names 
%of the products. We will introduce 3 types of vehicles, typeA can carry prodA, typeB can carry prodB and
%typeAB can carry both products. TypeAB has two sections with deterministic capacities. Each customer 
%has only one type of demand. We will have different variables for each vehicle type. We will solve three
%different pricing problems, one for each truck type in each facility. Pricing problem for typeA (typeB) includes
%only prodA (prodB) customers. All customers will exist in pricing problem for typeAB trucks. If we know the 
%exact solution of the LRS problem, this modified problem can be solved.       
%\end {itemize}
  
\section{Tentative Plan}
We plan to complete our future work based on the following schedule:

\begin{itemize}
\item We plan to complete the proposed algorithmic improvements by June 2006 to be able to solve
the LRS problem for instances with 50 customers.
\item We plan to spend 3 months to modify our formulation and algorithm to solve the special cases of the
LRS problem. We will be able to solve the LRPs, MDVRPs and VRPMTs by the end of September 2006. 
\item From October 2006 to May 2007, we will work on the stochastic LRS problem. 
\item We plan to work on the other extensions of the LRS problem after graduation. 
\end{itemize}    
